\documentclass[10pt,a4paper]{amsart}
\usepackage[margin=19mm]{geometry}
\usepackage{amsmath,amssymb,mathtools}
\usepackage[scr=rsfs,cal=euler]{mathalfa}
\usepackage{xfrac}
\usepackage[unicode,hidelinks,bookmarks=false]{hyperref}
\newcommand{\ie}{i.e.\ }
\newcommand{\cf}{cf.\ }
\newcommand{\resp}{respectively\ }
\newcommand{\Subsection}{\subsection}
\providecommand{\nobreakdash}{\nobreak\hbox{-}}
\setlength{\emergencystretch}{2em}
\multlinegap0pt
\usepackage{bm}
\usepackage{bbm}
\usepackage{dsfont}
\usepackage{xspace}
\usepackage{cancel}
\usepackage{mathtools}
\usepackage{cleveref}
\usepackage{amsthm}
\makeatletter
\@ifundefined{definition}{\newtheorem{definition}{Definition}}{}
\@ifundefined{lemma}{\newtheorem{lemma}{Lemma}}{}
\@ifundefined{proposition}{\newtheorem{proposition}{Proposition}}{}
\makeatother
\def\colim{\qopname\relax m{colim}}
\newcommand{\Con}{\mathrm{Con}}
\newcommand{\E}{\mathscr{E}}
\newcommand{\SSEG}{\mathbb{S}\mathrm{eg}}
\newcommand{\Seg}{\mathrm{Seg}}
\newcommand{\fun}{\mathrm{Fun}}
\newcommand{\im}{\mathrm{im}}
\newcommand{\op}{\mathrm{op}}
\title[Formal category theory in $\infty$-equipments I]{Formal category theory in $\infty$-equipments I}
\author{Jaco Ruit}
\date{Source: arXiv:2308.03583v3 (CC BY 4.0)}
\begin{document}
\maketitle

\noindent\emph{Source and licence.} Formal category theory in $\infty$-equipments I. Version arXiv:2308.03583v3 (CC BY 4.0), distributed under the Creative Commons Attribution 4.0 International licence, as recorded in the arXiv licence metadata for this exact version and shown on the abstract page https://arxiv.org/abs/2308.03583v3. Full rights notice: licenses/arxiv-2308.03583-CC-BY-4.0.txt.\par\smallskip

\noindent\textit{Editorial scope.} Counted unit: the Proposition whose source label is \texttt{prop.double-infty-cat-cat-objs} in the article (source0.tex lines 1595--1597, source label \texttt{prop.double-infty-cat-cat-objs}), delivered together with the article's own complete original proof of it (source0.tex lines 1598--1643, one \texttt{proof} environment of 46 lines). No mathematics is written, repaired or re-derived: statement and proof are the article's own text line for line. Required context retained uncounted: def.cat-modules (lines 1331--1347); prop.conduche-modules (lines 1373--1386); lemma.conduche-modules-kan-extension (lines 1459--1462). Every definition, display and label of those lines is retained unchanged. Omitted from this package and not redistributed: 1--1330, 1348--1372, 1387--1458, 1463--1594, 1644--3495. Nothing inside a retained excerpt is omitted or altered: every retained body is delivered exactly as the article's lines are written, so no omission receipt of any kind accompanies this package. Citations: the retained excerpts carry no citation command at all, so no bibliography entry of the article is transcribed and no reference is redistributed. Reference adaptation: none was needed and none was made; every reference inside the delivered unit resolves inside it, so no unresolved reference or citation is produced. Printed slips kept literal: no printed word, symbol, hypothesis or number of the counted statement or of its proof was altered. Layout and class: the article's own class is \texttt{amsart} with packages including amssymb, amsfonts, amsthm, kpfonts, eucal, geometry, tikz-cd, enumitem, hyperref, cleveref, aliascnt, todonotes, adjustbox, xstring; the wrapper is the lane's pinned \texttt{amsart} editorial wrapper at 10pt on A4, which restates the theorem-like environments the retained text needs and adds the article's own macro definitions that the retained text needs (\texttt{\textbackslash Con}, \texttt{\textbackslash E}, \texttt{\textbackslash SSEG}, \texttt{\textbackslash Seg}, \texttt{\textbackslash colim}, \texttt{\textbackslash fun}, \texttt{\textbackslash im}, \texttt{\textbackslash op}), with extra wrapper packages (bm, bbm, dsfont, xspace, cancel, mathtools, cleveref). The delivered numbering is the unit's own. Assets: no retained span contains a figure environment, a graphics-inclusion command, a PDF-page inclusion, a table or a TikZ picture; no image file is copied into this package, the package declares no asset dependency and no image file is redistributed with it. Licence: version arXiv:2308.03583v3 of this article is distributed under the Creative Commons Attribution 4.0 International licence, as recorded in the arXiv licence metadata for this exact version and shown on the abstract page https://arxiv.org/abs/2308.03583v3. A full rights notice is delivered at licenses/arxiv-2308.03583-CC-BY-4.0.txt. Characters: every retained excerpt is pure ASCII, so the delivered engine, XeLaTeX with its default Latin Modern text font, reports no missing character.
\begin{definition}\label{def.cat-modules}
	Let $n \geq 0$ be an integer.
	 A \textit{Segal $[n]$-module in $\E$} 
	is a functor 
	$
	M : (\Delta_{/[n]})^\op \rightarrow \E 
	$
	so that for any map $f : [k] \rightarrow [n]$ in $\Delta$, the comparison map
	$$
	M(f) \rightarrow M(f|\{0,1\}) \times_{M(f|\{1\})} \dotsb \times_{M(f|\{k-1\})} M(f|\{k-1,k\})
	$$
	is an equivalence. We will write 
	$$
	{\Delta_{/[n]}}\Seg(\E) \subset \fun((\Delta_{/[n]})^\op, \E)
	$$
	for the full subcategory spanned by the Segal $[n]$-modules.
\end{definition}

\begin{proposition}\label{prop.conduche-modules}
	Let $M$ be a Segal $[n]$-module in $\E$. Then, the following statements are equivalent: 
	\begin{enumerate}
		\item $M$ is left Kan extended along the inclusion $$(\Lambda_{/[n]})^\op \rightarrow (\Delta_{/[n]})^\op,$$
		\item for any map $\alpha : [2] \rightarrow [n]$, the restriction $\alpha^*M$ is left Kan extended along the inclusion 
		$$(\Lambda_{/[2]})^\op \rightarrow (\Delta_{/[2]})^\op,$$
		\item for any map $\alpha : [2] \rightarrow [n]$, the canonical map 
		$$
		\colim_{[k] \in \Delta^\op} M([k+2] \xrightarrow{0 \leq 1 \dotsb \leq 1\leq 2} [2] \xrightarrow{\alpha}
		 [n]) \rightarrow M([1] \xrightarrow{0\leq 2} [2] \xrightarrow{\alpha} [n])
		$$
		is an equivalence.
	\end{enumerate}
\end{proposition}

\begin{lemma}\label{lemma.conduche-modules-kan-extension}
	Suppose that $N : (\Lambda_{/[n]})^\op \rightarrow \E$ satisfies the Segal condition as in \cref{def.cat-modules}. Then 
	the left Kan extension $M : (\Delta_{/[n]})^\op \rightarrow \E$ of $N$ along $(\Lambda_{/[n]})^\op \rightarrow (\Delta_{/[n]})^\op$ satisfies the Segal condition as well.
\end{lemma}

\begin{proposition}\label{prop.double-infty-cat-cat-objs}
	The simplical $\infty$-category $\SSEG(\E)$ is a double $\infty$-category.
\end{proposition}

\begin{proof}
	Let $\Delta_{/[n]}\Con'(\E) \subset \fun((\Lambda_{/[n]})^\op, \E)$ be the full subcategory spanned 
	by those functors that satisfy the Segal condition as in \cref{def.cat-modules}.
	Then it readily 
	follows from characterization (1) of \cref{prop.conduche-modules} and \cref{lemma.conduche-modules-kan-extension} that the fully faithful left Kan extension functor 
	$
	i_!: \fun((\Lambda_{/[n]})^\op, \E) \rightarrow \fun((\Delta_{/[n]})^\op, \E)
	$
	along $i : (\Lambda_{/[n]})^\op \rightarrow (\Delta_{/[n]})^\op$
	restricts to an equivalence between $\Delta_{/[n]}\Con'(\E)$ and $\Delta_{/[n]}\Con(\E)$. Consequently, 
	we see that it suffices to check that the canonical functor 
	$$
	p : \Delta_{/[n]}\Con'(\E) \rightarrow {\Delta_{/[1]}}\Seg(\E) \times_{\Seg(\E)} \dotsb \times_{\Seg(\E)} {\Delta_{/[1]}}\Seg(\E)
	$$
	is an equivalence.

	To this end, we note that this functor may be identified with the restriction of the left adjoint 
	that sits in the adjunction
	$$
	j^* : \fun((\Lambda_{/[n]})^\op, \E) \rightleftarrows \fun((\Lambda_{/[n]})^{\mathrm{cell}, \op},\E)  : j_*
	$$
	associated to \textit{right} Kan extending along the inclusion
	$$j : (\Lambda_{/[n]})^{\mathrm{cell},\op} \simeq (\Delta_{/[1]})^\op \cup_{\Delta^\op} \dotsb \cup_{\Delta^\op} (\Delta_{/[1]})^\op \rightarrow (\Lambda_{/[n]})^\op$$
	of the full subcategory of those cellular maps $f : [m] \rightarrow [n]$ so that $\im(f) \subset \{i,i+1\}$ for some $i$. 
	Let us write $\Delta_{/[n]}\Con''(\E) \subset \fun((\Lambda_{/[n]})^{\mathrm{cell}, \op},\E)$ for the full subcategory of 
	those functors that satisfy the Segal condition. Then $p$ may be identified with the restriction of $j^*$ 
	to $\Delta_{/[n]}\Con'(\E)$ and $\Delta_{/[n]}\Con''(\E)$. 
	Since $j_*$ is fully faithful, we thus see that $p$ is an equivalence if the following holds:
	\begin{enumerate}
		\item the left adjoint $j_*$ carries $\Delta_{/[n]}\Con''(\E)$ into $\Delta_{/[n]}\Con'(\E)$,
		\item if $M \in \Delta_{/[n]}\Con'(\E)$, then $M$ is right Kan extended along $j$, i.e.\, the unit
		$M \rightarrow j_*j^*M$ is an equivalence,
	\end{enumerate}
	in which case $j_*$ restricts to an inverse of $p$. 
	
	Both (1) and (2) follow from the following computation. Suppose that 
	 $M : (\Lambda_{/[n]})^{\mathrm{cell}, \op} \rightarrow \E$ is a functor, and let $f : [m] \rightarrow [n]$ be a cellular map, then 
	$$
	j_*M(f) \simeq M(f|f^{-1}(\{\alpha, \alpha+1\})) \times_{M(f|f^{-1}\{\alpha+1\})} \dotsb \times_{M(f|f^{-1}\{\omega-1\})} M(f|f^{-1}(\{\omega-1,\omega\})),
	$$
	where $\alpha := \min(f)$ and $\omega := \max(f)$. 
	Consequently, if $M$ satisfies the Segal condition then we get that
	\begin{displaymath}
	j_*M(f) \simeq M(f|\{0,1\}) \times_{M(f|\{1\})} \dotsb \times_{M(f|\{m-1\})} M(f|\{m-1,m\}). \qedhere
	\end{displaymath}
\end{proof}

\end{document}
